\section{The Evaluation Rubric}
\label{sec:rubric}

The rubric is the central contribution of this survey
(Contribution C1, \S\ref{sec:intro:contributions}) and the
primary instrument by which subsequent sections compare the
eight protocols of the tool-binding and peer-messaging
layers. It comprises ten evaluation axes spanning
discovery, identity, authorization, transport, task
semantics, async behaviour, multi-modal handling, security,
semantic interoperability, and governance. Seven of the ten
axes (Axes 1, 2, 3, 5, 6, 7, and 9) use anchored ordinal
scoring on a 0 to 3 scale; one (Axis 4, Transport) is a
structured tuple; one (Axis 8, Security) is a count over a
fixed twelve-item threat catalogue; one (Axis 10,
Governance/License/Observability) is a composite triple
combining a categorical governance class, a license
identifier, and a 0 to 3 observability score. The asymmetry
between ordinal, tuple, count, and composite forms is
deliberate (\S\ref{sec:rubric:goals}): forcing every axis
into the same numerical form would impose a false precision
on dimensions where it is not warranted.

\subsection{Design goals}
\label{sec:rubric:goals}

Four design constraints shaped the rubric.

\paragraph{D1: Uniform applicability across heterogeneous
protocols.} The rubric must apply consistently to a
tool-binding protocol (MCP) and a peer-messaging protocol
(A2A) and a commerce-settlement protocol (AP2) without
re-defining axes per layer. Where an axis is structurally
inapplicable to a given protocol (e.g., multi-modal content
handling for a settlement-attestation protocol), the score
records \emph{not applicable} rather than artificially low.

\paragraph{D2: Anchored descriptors, not free-form numerals.}
Each ordinal level (0, 1, 2, 3) for an ordinal axis carries an
explicitly named anchor descriptor in this section. A score
of 2 for Axis 1 (Discovery) on protocol P means \emph{P
exhibits the level-2 descriptor}, not \emph{P is somewhere
between protocols Q and R}. Anchoring is the only way to
prevent the score from collapsing into the surveyor's
subjective ordering, which is the principal weakness of the
qualitative comparison tables in the prior
surveys~\cite{yang2025survey, ehtesham2025interop}.

\paragraph{D3: Refusal to manufacture composite scores.}
The rubric is presented as a multi-dimensional matrix
(Table~\ref{tab:rubric-matrix}), a heatmap
(Figure~\ref{fig:heatmap}), and per-protocol radar charts
(Figure~\ref{fig:radar}). No single composite score is
computed. Composite scores would invite ranking, ranking
would invite optimization-by-vendor, and the resulting
benchmark-gaming dynamic has repeatedly distorted
comparable agent-evaluation work. Per-axis scores remain
defensible; their sum or weighted average does not.

\paragraph{D4: Re-runnable figures from a machine-readable
source of truth.} Every score in Table~\ref{tab:rubric-matrix},
every cell in Figure~\ref{fig:heatmap}, and every spoke in
Figure~\ref{fig:radar} is regenerated from
\texttt{data/protocol-scores.yaml} by the scripts in the
companion repository (Contribution C4,
\S\ref{sec:intro:contributions}). A future reader who
disputes a score can fork the YAML, change the cell, and
re-run the figure generation; the survey's claims are
re-runnable rather than only re-readable. This goal
concerns the path from a score to a figure. It says nothing
about whether a second rater would assign the same score,
which is a separate property that
\S\ref{sec:rubric:calibration} addresses on its own terms.

\subsection{The ten axes}
\label{sec:rubric:axes}

\subsubsection{Axis 1: Discovery and Registry Model}
\label{sec:rubric:axis1}

This axis evaluates how an agent locates another agent and
acquires the metadata necessary to interoperate. It captures
both the discovery primitive (well-known endpoint, DNS-like
directory, DID resolver, content-addressed registry) and the
strength of the verifiability guarantees the discovery layer
provides.

\begin{description}
  \item[Score 0: Out-of-band only.] No protocol-level
    discovery primitive. Agents must be configured with
    endpoints through external means (configuration files,
    environment variables, hand-edited registries). MCP in
    its stdio-only configuration sits near this baseline.
  \item[Score 1: Well-known endpoint or
    config-mediated.] A documented convention by which a
    client can probe a host for a service descriptor
    (e.g., A2A's \texttt{/.well-known/agent-card.json}~\cite{spec_a2a}).
    Centralized; trust delegated to the hosting authority.
  \item[Score 2: Federated directory with structured
    descriptors.] A directory-like service (registry,
    federated catalogue, on-chain index) returning typed
    descriptors. ACNBP's Agent Name Service~\cite{spec_acnbp}
    and Coral's MCP-mediated \texttt{list\_agents} flow~\cite{spec_coral}
    sit here.
  \item[Score 3: Decentralized verifiable resolution.]
    Discovery returns cryptographically-verifiable
    descriptors resolvable independently of any single
    authority. Typified by DID-resolver-based discovery in
    ANP~\cite{spec_anp} and LOKA's Semantic Discovery
    Fabric~\cite{spec_loka}.
\end{description}

Evidence supporting an Axis 1 score must point to a specific
section of the protocol specification that describes the
discovery primitive, the descriptor schema, and the
verification mechanism (or its explicit absence).

\subsubsection{Axis 2: Identity and Authentication}
\label{sec:rubric:axis2}

This axis evaluates the mechanism by which an agent
establishes who it is to its counterparties and how that
claim can be verified. Modern inter-agent protocols span a
range from process-bound identity (the stdio-invoked
process inherits its parent's trust) through bearer-token
schemes (OAuth~\cite{rfc6749}) to fully decentralized
identifiers (DIDs~\cite{w3c_did_core}) with verifiable
credentials~\cite{w3c_vc_data_model}. The anchors below name
those three standards, so a reader can check a score against
the normative text rather than against our paraphrase.

\begin{description}
  \item[Score 0: Process trust or unspecified.] No
    protocol-level identity mechanism. Trust derives from
    process-control (the parent invoked the child) or is
    deferred to the deployment environment.
  \item[Score 1: Bearer tokens (OAuth 2.x or
    equivalent).] Identity carried by a token issued by a
    third-party authority. A2A in its standard
    configuration~\cite{spec_a2a}, MCP over HTTP+OAuth
    2.1~\cite{spec_mcp}, and ACP all sit at level 1.
  \item[Score 2: Bearer tokens hardened with mTLS or
    signed envelopes.] Bearer credentials augmented with
    cryptographic transport binding or message-level
    signatures resistant to bearer-token theft. ACNBP's
    X.509 + OCSP + ECC-signed envelope
    structure~\cite{spec_acnbp} sits at level 2.
  \item[Score 3: Decentralized identifiers with
    verifiable credentials.] Identity established by a
    cryptographic identifier (typically a
    DID~\cite{w3c_did_core}) whose
    associated keys can sign arbitrary
    challenge-response exchanges, paired with W3C
    Verifiable Credentials~\cite{w3c_vc_data_model} for
    attribute assertions.
    ANP~\cite{spec_anp}, LOKA~\cite{spec_loka}, and Coral
    (DIDs plus Solana wallets) sit at level 3.
\end{description}

The prior survey by Kong et al.~\cite{kong2025communication}
treats identity primarily through the security lens; we
score it separately here to distinguish the
\emph{identity-modeling} contribution of a specification
from its \emph{threat-mitigation coverage} (Axis 8).

\subsubsection{Axis 3: Authorization and Delegation Model}
\label{sec:rubric:axis3}

This axis evaluates how a protocol expresses what an
authenticated agent is allowed to do, and how authorization
propagates across delegation chains (agent A asks agent B
to do X on behalf of principal P). The contemporary
literature treats this as the most under-developed area of
the contemporary protocols~\cite{authz_propagation,
agentic_jwt}; the rubric captures the gap explicitly.

\begin{description}
  \item[Score 0: Unspecified.] No protocol-level
    authorization primitives. Authorization, if any, is
    enforced at the application layer.
  \item[Score 1: Scope tokens or token forwarding.]
    OAuth-style scopes carried in bearer tokens; or the
    token-forwarding pattern in which agent A passes its
    own access token onward to agent B. Carries no
    cryptographic delegation chain; the receiver cannot
    distinguish ``the principal authorized this agent''
    from ``a previous agent in the chain authorized this
    onward call.'' A2A sits at level 1 because it
    standardises no delegation credential: \S7.5 of the
    specification makes authorization logic
    implementation-specific and leaves the choice of
    delegation mechanism to the surrounding identity
    architecture~\cite{spec_a2a}. Deployments then adopt
    token forwarding, which is where the security
    literature observes it~\cite{authz_propagation,
    agentic_jwt}. MCP's OAuth scopes also sit at level 1.
    The level is an omission at the specification layer,
    not a pattern either specification prescribes.
  \item[Score 2: Verifiable-credential-scoped
    capabilities or structured delegation tokens.]
    Authorization expressed as a signed capability with
    explicit scope, audience, and time bounds, but
    without a fully standardized chain-of-custody format.
    ANP's VC-scoped capabilities~\cite{spec_anp}, LOKA's
    smart-contract lifecycle policies~\cite{spec_loka}, and
    ACNBP's Capability Binding ceremony~\cite{spec_acnbp}
    sit at level 2. Each binds authority cryptographically
    between two named parties. None specifies how that
    authority travels further.
  \item[Score 3: Cryptographically-bound capability
    chains.] Multi-step delegation expressed as a
    composable chain of signed capabilities (or a
    capability-binding ceremony) in which each
    intermediate agent's authorization is independently
    verifiable. The second clause is the demanding one: a
    downstream agent must be able to check an upstream
    grant it did not witness. No peer-messaging protocol in
    the scored set reaches this level
    (\S\ref{sec:comparison:derivation}). The
    commerce-settlement protocols
    (\S\ref{sec:industry-protocols}) do reach it, by
    construction, through their intent-mandate-receipt
    envelope, and the agentic-JWT
    proposal~\cite{agentic_jwt} is the leading candidate
    for porting the level to the peer-messaging layer.
\end{description}

\subsubsection{Axis 4: Transport and Wire Format}
\label{sec:rubric:axis4}

The transport axis is recorded as a structured tuple rather
than an ordinal score. A single number cannot capture the
joint specification of transport bindings, serialization
format, streaming mechanism, and async-delivery primitive.
The tuple has the following fields, each populated from the
protocol specification:

\begin{itemize}
  \item \texttt{transports}: ordered list of binding
    technologies (e.g., \texttt{HTTPS/JSON-RPC 2.0},
    \texttt{stdio}, \texttt{HTTP+SSE}, \texttt{streamable-HTTP},
    \texttt{gRPC}, \texttt{REST}, \texttt{WebSocket}).
  \item \texttt{serialization}: payload encoding
    (e.g., \texttt{JSON}, \texttt{JSON-LD}, \texttt{multipart},
    \texttt{Protocol Buffers}).
  \item \texttt{streaming}: in-protocol streaming mechanism
    if any (\texttt{SSE}, \texttt{chunked HTTP},
    \texttt{WebSocket frames}, \texttt{none}).
  \item \texttt{async}: long-running-operation delivery
    primitive (\texttt{push notifications},
    \texttt{polling}, \texttt{webhooks},
    \texttt{durable resumable tasks}, \texttt{intent-centric
    peer-to-peer}).
\end{itemize}

The tuple is reported alongside the ordinal axes in
Table~\ref{tab:rubric-matrix}; it is not aggregated. Its
purpose is to support the comparison of like-with-like
transport choices (e.g., MCP's \texttt{stdio} versus
\texttt{streamable-HTTP} variants) without collapsing
distinct engineering trade-offs into a single
``transport richness'' score.

\subsubsection{Axis 5: Task and Conversation Model}
\label{sec:rubric:axis5}

This axis evaluates whether the protocol treats
multi-message interactions as first-class objects with a
defined lifecycle, or treats them as ad-hoc sequences of
independent request-response exchanges.

\begin{description}
  \item[Score 0: Stateless request/response only.] No
    notion of a multi-message interaction; each message
    stands alone. The receiver carries no protocol-level
    state from one call to the next.
  \item[Score 1: Session-scoped but no task object.]
    A session identifier groups related messages but the
    protocol defines no terminal status or artifact
    aggregation. MCP in its standard request/response
    configuration sits near level 1.
  \item[Score 2: Persistent multi-participant
    interaction.] A stateful conversation object groups
    messages, supports multiple participants, but lacks a
    full lifecycle state machine with sub-task
    delegation. Coral's persistent threads~\cite{spec_coral}
    sit at level 2.
  \item[Score 3: Full task object with lifecycle,
    artifacts, and sub-task delegation.] Tasks are
    first-class with defined lifecycle states (e.g.,
    \emph{submitted}, \emph{working}, \emph{input-required},
    \emph{completed}, \emph{failed}, \emph{canceled}), an
    artifact aggregation model, and primitives for
    delegating to a sub-task. A2A~\cite{spec_a2a} and
    ACNBP's 10-step bound-execution
    machine~\cite{spec_acnbp} sit at level 3.
\end{description}

\subsubsection{Axis 6: Async and Long-Running Operations}
\label{sec:rubric:axis6}

This axis evaluates the protocol's support for operations
whose duration exceeds the lifetime of a single connection
or even the lifetime of either communicating agent. It is
distinct from Axis 5 (which evaluates whether tasks exist
as first-class objects) because a protocol can offer rich
task semantics over short-lived connections and still
struggle with multi-hour operations.

\begin{description}
  \item[Score 0: No async primitive.] All operations
    complete within a single request/response cycle or
    fail.
  \item[Score 1: Polling-based.] Long-running
    operations are tracked by an identifier that the
    client must poll. No push delivery from server to
    client.
  \item[Score 2: Server-initiated notifications.]
    The server can push notifications or events back to
    the client via a long-lived stream (SSE,
    WebSockets) but the conversation is bound to the
    open connection.
  \item[Score 3: Durable resumable tasks.] Operations
    survive connection loss; on reconnect, the client
    resumes the task by its identifier and receives a
    complete or partial result and any backfilled
    events. A2A's push notifications combined with
    durable resumable tasks~\cite{spec_a2a} sits at
    level 3; LOKA's checkpoint-and-recovery
    machinery~\cite{spec_loka} is conceptually at level
    3 though implementation-incomplete.
\end{description}

\subsubsection{Axis 7: Multi-Modal Content Handling}
\label{sec:rubric:axis7}

This axis evaluates the protocol's vocabulary for content
beyond plain text. Multi-modal handling has emerged as a
first-class concern only in the 2024 to 2026 wave; older
ACLs and most service-oriented stacks treated content as
opaque payload with at most a MIME type.

\begin{description}
  \item[Score 0: Plain text payload only.] Messages
    carry text; binary or structured content is the
    caller's problem.
  \item[Score 1: Text plus referenced attachments.]
    Messages can reference external content via URLs or
    file handles but do not specify a typed inline part
    model.
  \item[Score 2: Typed content blocks.] A small set of
    typed parts (text, image, structured data, file)
    that can be composed within a single message. MCP's
    typed content blocks~\cite{spec_mcp} sit at level 2.
  \item[Score 3: Typed parts with inline-or-reference
    duality and content negotiation.] A complete part
    model in which each part declares its MIME type, its
    transmission mode (inline or by reference), and can
    participate in capability negotiation between
    sender and receiver. A2A's TextPart/FilePart/DataPart
    model~\cite{spec_a2a} sits at level 3.
\end{description}

\subsubsection{Axis 8: Security Threat Surface}
\label{sec:rubric:axis8}

This axis is the only count-based axis in the rubric. We
maintain a fixed twelve-item catalogue of inter-agent
protocol threats, drawn from synthesis of the security
literature~\cite{kong2025communication, a2asecbench,
zhan2024injecagent, agentleak, protocol_exploits,
a2a_weaknesses, threatmodel2026, did_vc_agents} and
applied uniformly:

\begin{enumerate}
  \item Authentication bypass
  \item Replay attacks
  \item Message tampering / man-in-the-middle
  \item Supply-chain attacks (registry poisoning,
    descriptor tampering)
  \item Denial-of-service
  \item Side-channel privacy leakage
  \item Decentralized-or-untrusted registry compromise
  \item Capability escalation through ambiguous
    descriptors
  \item Cross-agent privilege confusion
  \item Foundation-model layer prompt injection
  \item Audit-trail tampering / non-repudiation breach
  \item Delegation-chain forgery
\end{enumerate}

A protocol's Axis 8 score is the count $N \in [0, 12]$ of
threats explicitly addressed at the specification level
(not at the deployment level), with citations to the
specification sections that establish each mitigation. The
threat-catalogue and the per-protocol coverage details
are presented as the threat-coverage matrix
in \S\ref{sec:security}. The axis is intentionally a
\emph{count} rather than an ordinal score because the
twelve threats are heterogeneous and weighting them inside
the axis would re-introduce subjective ranking. A high
count is a necessary but not sufficient condition for
protocol security; \S\ref{sec:security} discusses the
difference.

\paragraph{Why the count alone is not enough, and what we
report beside it.}
A count treats every threat as one unit. Two protocols can
score the same and leave very different amounts of risk
behind, because one of them addressed the two cheapest
classes. Registry poisoning hands an attacker every
consumer of a poisoned descriptor; denial of service costs
availability and recovers when the load stops. Counting
them equally is the weakness a reviewer of this rubric
identified, and the count cannot answer it on its own.

Our answer is a second view, not a replacement score.
Table~\ref{tab:severity-coverage} in \S\ref{sec:security}
reports \emph{severity-weighted coverage} beside the count.
Each of the twelve classes carries a band (critical, high,
or moderate) and an integer weight (3, 2, or 1). The table
reports the addressed weight over the applicable weight,
and the unaddressed weight in the critical and high bands.

Three stated criteria assign the band. Does exploitation
let the attacker make the agent \emph{act}, rather than
only observe? Is the compromise recoverable from evidence
the protocol itself keeps? Does one exploit reach beyond
the pair of agents involved? Those criteria borrow their
structure from the OWASP AI Vulnerability Scoring
System~\cite{owasp_aivss}. That work argues that agentic
risk is dominated by loss of control over agent actions,
not over data alone. We do not report an AIVSS score.
AIVSS scores a concrete vulnerability in a concrete
deployment. We band a threat class against a
specification, which carries no exploit-maturity signal
and no environment.

This is deliberately consistent with D3 rather than an
exception to it. The severity view is a second reported
column, not a collapsed protocol quality score. Its
weights are published in the companion repository, so a
reader who disputes a band can change the weight and
regenerate the table. What D3 refuses is a single number
standing in for ten axes. Reporting a count and a
severity-weighted view of the same axis, with the weights
exposed, is the opposite move. It makes the weighting
argument explicit instead of hiding it inside an ordinal.

\subsubsection{Axis 9: Semantic Interoperability}
\label{sec:rubric:axis9}

This axis evaluates the protocol's mechanism for ensuring
that the meaning of an exchanged capability description,
intent, or parameter is interpretable by parties that did
not jointly author the descriptor. Section~\ref{sec:semantic}
treats this as a distinct concern from message passing.

\begin{description}
  \item[Score 0: Free-text only.] Descriptors are
    natural-language strings with no machine-tractable
    semantics.
  \item[Score 1: Structured schemas with free-text
    descriptions.] Typed input/output schemas (e.g.,
    JSON Schema) for parameters, but the human-meaningful
    description of what a capability does remains free
    text. MCP, A2A, and Coral all sit at level 1.
  \item[Score 2: Ontology-base or semantic-discovery
    indexing.] An ontology, vocabulary, or
    semantic-indexing mechanism that constrains
    capability descriptors to a shared term space.
    ANP's JSON-LD vocabulary with Agent Description
    Protocol~\cite{spec_anp} sits at level 2 at the
    cutoff; \S\ref{sec:semantic:ontology} records a
    post-cutoff narrowing of that commitment.
  \item[Score 3: Full capability ontology with
    semantic matching.] A specified ontology of
    capabilities and intents, supported by a matching
    mechanism (semantic-similarity search or
    logic-based reasoning) that produces verifiable
    semantic alignments. LOKA's Polyglot Intent Engine
    and capability ontologies~\cite{spec_loka} are
    aspirational level 3; no surveyed protocol
    \emph{ships} level 3 as of 2026-Q1.
\end{description}

The recurring gap between level 1 (where the LLM-era
protocols cluster) and level 3 (where FIPA-SL and full
ontology commitments once aimed) is the principal subject
of \S\ref{sec:semantic} and of the open-problems
agenda (\S\ref{sec:openproblems}).

\subsubsection{Axis 10: Governance, License, and Observability}
\label{sec:rubric:axis10}

The final axis is a composite triple recording three
institutional and operational characteristics that the
prior surveys treat only in passing:

\begin{itemize}
  \item \textbf{Governance}: the legal entity or
    community that holds standardization authority over
    the specification. We record the entity by name
    (e.g., \emph{Linux Foundation}, \emph{Linux
    Foundation AAIF}~\cite{lf_aaif},
    \emph{community / W3C-aligned}~\cite{w3c_agent_protocol_cg},
    \emph{individual authors}, \emph{Coral Protocol team},
    \emph{academic authors}).
  \item \textbf{License}: the license under which the
    specification (and any reference implementation) is
    distributed (e.g., Apache~2.0, MIT-equivalent,
    CC-BY-4.0, CC-BY-NC-SA~4.0).
  \item \textbf{Observability}: an ordinal 0 to 3 score on
    the protocol's specification of audit, tracing, and
    monitoring primitives (level 0 = unspecified;
    level 1 = logging hooks emerging; level 2 = trace
    fields specified, no audit primitive; level 3 =
    immutable audit trail and observability primitives
    specified).
\end{itemize}

The triple of steward, license, and observability score
is reported as a single composite cell in
Table~\ref{tab:rubric-matrix} and is unpacked into its
component analyses in \S\ref{sec:governance}.

\subsection{Scoring methodology}
\label{sec:rubric:method}

Scores are derived from a structured reading of each
protocol specification at its declared surveyed version
(2026-Q1 cutoff). The procedure for each
\emph{(protocol, axis)} pair is:

\begin{enumerate}
  \item Identify the specification section(s) relevant to
    the axis.
  \item Map the section's behaviour to the anchored
    descriptors of the axis.
  \item Record the resulting score, a one-sentence
    evidence summary, and the citation(s) supporting the
    evidence.
  \item Where a protocol specification is ambiguous, the
    score reflects what the specification \emph{requires},
    not what reference implementations \emph{do}. A
    facility provided only by a non-canonical
    implementation is recorded in evidence but does not
    elevate the score.
\end{enumerate}

Scores are stored as \texttt{data/protocol-scores.yaml} in
the companion repository (Contribution C4) and regenerate
Table~\ref{tab:rubric-matrix} and Figures~\ref{fig:heatmap}
and~\ref{fig:radar} mechanically. Anyone more interested
in the basis of a given score should consult the cited
specification section first and the YAML entry second; the
entry will name the specific evidence the score rests on.

\paragraph{What the rubric does not produce.} The rubric
does not produce a single composite ``protocol quality''
score. It does not produce a ranking. It does not produce
a recommendation: the appropriate choice of protocol
depends on deployment-specific axes that lie outside the
specification (operational maturity, ecosystem support,
internal toolchain compatibility) and that no
spec-reading exercise can determine. The rubric produces
a \emph{normalized comparison surface}; the user of the
comparison brings the application-specific weighting.

\subsection{Calibration, adjudication, and rater agreement}
\label{sec:rubric:calibration}

The procedure above says how a score is derived. It does
not say how two people applying it reach the same number.
An anchored rubric that different readers score differently
is not yet a normalized comparison surface. This subsection
closes that gap in three parts. First, the calibration
rules a rater applies at a contested boundary. Second, two
worked examples of those rules deciding a real case. Third,
an honest statement of what agreement evidence this survey
does and does not carry.

\paragraph{Four calibration rules.}
Every contested call in this survey was settled by the same
four rules, applied in order.

\begin{enumerate}
  \item \textbf{Normative text beats prose.} A mechanism
    counts at the level the specification \emph{requires}
    it. Language in an overview, a rationale, or a
    non-normative appendix supports at most the level
    below.
  \item \textbf{Spec beats implementation.} A facility that
    only a reference implementation provides is recorded in
    the evidence field and does not raise the score. This
    is rule 4 of \S\ref{sec:rubric:method}, restated here
    because it decides more boundary cases than any other.
  \item \textbf{The lower anchor wins a tie.} If a rater
    can argue the case at either of two adjacent levels,
    the score is the lower one. The asymmetry is
    deliberate. It makes the rubric under-credit rather
    than over-credit, which is the safer error for a
    reader deciding whether to depend on a protocol.
  \item \textbf{Absence of a primitive is not
    inapplicability.} \emph{Not applicable} is reserved for
    a threat or facility that is structurally outside the
    protocol's scope, not for one the protocol simply
    omits. MCP has no AgentCard, so the AgentCard
    supply-chain class is \texttt{na} for MCP. ACP has an
    unsigned manifest that plays the same role, so the same
    class is \texttt{not\_addressed} for ACP, not
    \texttt{na}.
\end{enumerate}

\paragraph{Worked example 1: ACP, Axis 2 and the
impersonation cell.}
An early pass scored ACP's identity layer as
OAuth-scope-based, on the strength of the specification's
prose about deployment behind an authenticating reverse
proxy. Re-reading the OpenAPI document showed it declares
no \texttt{securitySchemes} at all. Rule 1 settles it: the
prose is not normative, the machine-readable contract is,
and a mechanism the contract does not declare cannot be
scored. Rule 2 blocks the fallback argument that ACP
deployments do authenticate in practice. The impersonation
cell moved to \texttt{not\_addressed} and ACP's Axis 8
count moved to 1. A second rater who applies rules 1 and 2
to the same OpenAPI document reaches the same cell. The
question the rules ask has a yes-or-no answer in the
document rather than a judgement.

\paragraph{Worked example 2: AGNTCY, the audit-log
tampering cell.}
AGNTCY persists scan reports as OCI referrers and emits
OpenTelemetry traces. Both are real, specified, and
relevant to audit. Neither is a tamper-evident log: nothing
in the specification prevents a party with write access
from rewriting the record. The candidate scores are
\texttt{addressed} (an audit mechanism is specified) and
\texttt{partial} (the mechanism is specified but the
integrity property is not). Rule 3 decides it as
\texttt{partial}, and the cell's evidence field records the
reason in the same words. The same rule applied to Coral's
blockchain-anchored trail returns \texttt{addressed},
because there the integrity property is the mechanism. The
contrast is what makes the boundary reproducible: the rater
is not asked whether the audit story is good, but whether
the specification names an integrity guarantee.

\paragraph{Adjudication when the rules do not settle it.}
Two cases survive the four rules. The first is a
specification whose text is genuinely ambiguous about what
it requires. Here we wrote to the specification's authors,
and where they answered, the answer is recorded as evidence
and the score follows it. Two scores in this survey moved
that way. ACNBP's Axis 3 dropped from 3 to 2 after a lead
author confirmed that the Capability Binding ceremony is a
two-party protocol. Multi-agent delegation is outside its
scope (\S\ref{sec:security:authz}). ANP's task-model scope
was recorded as deliberate rather than underspecified after
the same kind of exchange. The second
case is a specification that changed after we scored it.
There the cutoff governs: the score stands at the surveyed
version, and the drift is recorded rather than silently
absorbed (\S\ref{sec:rubric:limits}).

\paragraph{What agreement evidence we have, and what we do
not.}
We state this plainly, because the rubric's value depends
on it. This survey does not report an inter-rater
reliability statistic. No second rater independently scored
the full matrix. We therefore have no Cohen's $\kappa$ and
no agreement percentage to offer, and we decline to compute
one from a process that would not support it.

What we have instead is an auditable change history. Every
score lives in \texttt{protocol-scores.yaml} with an
evidence string. Every change to a score is a commit with a
stated reason. The counting rule for Axis 8 is enforced by
a check in the companion repository, which fails the build
when a count and the threat matrix disagree. That check
exists because they did disagree: four protocols carried
counts that had absorbed \texttt{partial} cells, which the
axis definition excludes, and nothing caught it because
nothing was checking. Coverage of the manuscript prose was
added to the same check for the same reason.

An auditable trail is weaker evidence than two independent
raters, and we do not present it as equivalent. It supports
a narrower claim: that each score is traceable to a
specification passage a disputing reader can check, and
that no score has moved without a recorded reason. The
axes most exposed by the gap are the two most
judgement-heavy ones: semantic interoperability (Axis 9)
and the observability component of Axis 10. In both, the
anchors describe a degree of specification rather than the
presence of a named mechanism. A two-rater pass over them
is the highest-value item in the rubric's own future work,
and \S\ref{sec:openproblems:benchmarking} lists it there.

\subsection{Limitations of the rubric}
\label{sec:rubric:limits}

Four limitations of the rubric require explicit
acknowledgement.

\paragraph{The rubric covers eight of the thirteen surveyed
protocols.} Every score in this survey is a score of one of
the eight tool-binding or peer-messaging protocols. The five
commerce-settlement protocols (AP2, Visa TAP, Mastercard
Agent Pay, ERC-8004, UCP) are surveyed and compared, in
\S\ref{sec:emerging:commerce} and
Table~\ref{tab:commerce-stack}, but they are not rubric
scored. The reason is that six of the ten axes have no
referent in a settlement envelope: it has no task state
machine, no async resumption model, no multi-modal content
negotiation, and its discovery and transport are those of
whichever peer-messaging protocol carries it. Scoring them
anyway would produce rows of \texttt{n/a} and a set of
figures whose averages moved for reasons of scope rather
than of substance.

UCP is the one member of the layer where that reason is
weaker. It specifies discovery, authentication,
authorization, and four transport bindings
(\S\ref{sec:emerging:commerce}), so four axes would return a
defensible score. It is left unscored anyway, on two
grounds. It is a commerce-orchestration layer rather than a
peer-messaging protocol, so the axes it does answer are
answered for a different purpose. And scoring one member of
the commerce layer while four remain unscored would move the
figures' averages for reasons of scope. Its per-axis
evidence is recorded in
\texttt{data/protocol-scores.yaml} so that a later revision
widening the rubric does not start from nothing. The cost of
the choice is that the commerce layer cannot be ranked
against the other two, and a reader wanting one comparable
number across all thirteen will not find it here.

\paragraph{Subjectivity at the anchor boundaries.} The
boundary between two adjacent anchored levels is
necessarily a judgment call. We have made each boundary
description as concrete as the specification language
allows, and the four calibration rules of
\S\ref{sec:rubric:calibration} settle a contested boundary
in a stated order rather than by preference. Still, some
boundaries remain thin. A \emph{level-1} discovery primitive
is a well-known endpoint and a \emph{level-2} one is a
federated directory. Which side A2A's AgentCard convention
falls on depends on whether one reads it as a
``federation'' or as ``a
hosted convention.'' Where the boundary is contestable, we
record the reasoning in the YAML evidence field. The
residual limitation is the one
\S\ref{sec:rubric:calibration} states directly: the rules
make a boundary call reproducible in principle, and no
second rater has independently confirmed that it
reproduces.

\paragraph{Version drift across the review cycle.}
Protocols at the active 2025 to 2026 frontier are evolving
on a quarterly cadence. A score recorded against
A2A v1.0 may not hold against the version current at
publication. The companion repository's update discipline
addresses this directly: the YAML carries
\texttt{spec\_version} and \texttt{retrieved} fields per
protocol; a future reader can compute the staleness
delta. The survey itself reports scores at the declared
cutoff and does not silently update them in proof.

\paragraph{Evidence tier asymmetry.} Some protocols
(A2A, MCP) have multi-tier sourcing: an official
specification, peer-reviewed security analyses, and
production deployments. Others (ACNBP, LOKA) are at
arXiv-paper stage with no formal SDO and no large-scale
production deployment to validate the spec's
operational claims. Where the specification declares a
mechanism the rubric scores the declaration; where the
declaration depends on infrastructure that does not
yet exist (LOKA's level-3 semantic interoperability via
the Polyglot Intent Engine), the evidence field records
the gap. We are weighing scores by source tier;
\S\ref{sec:industry-protocols} and
\S\ref{sec:emerging-protocols} present the
implementation-maturity context alongside the rubric
scores.
